\exercisesection{A-模复形}

\begin{enumerate}
\def\labelenumi{\arabic{enumi})}
\tightlist
\item
  假设环 A 是主理想环；设 C 是自由 A-模的一个复形。
\end{enumerate}

\begin{enumerate}
\def\labelenumi{\alph{enumi})}
\item
  证明 C 是一族复形 \((^{p}\mathbf{E})_{p\in \mathbb{Z}}\) 的直和，满足： \(^{p}\mathbf{E}_{n} = 0\) 当 \(n\neq p\) , \(p - 1 ; Z_{n}(^{p}\mathbf{E}) = 0\) 。
\item
  我们进一步假设 \(C_n\) 对所有 \(n \in \mathbf{Z}\) 都是有限生成的。证明 \(C\) 是一些复形的直和，这些复形仅有一个非零分量或两个连续的非零分量，这些分量都同构于 \(A\) 。
\end{enumerate}

\begin{enumerate}
\def\labelenumi{\arabic{enumi})}
\setcounter{enumi}{1}
\item
  设 C 与 C' 为两个 A-模复形；假设复形 C 与 B(C)（相应地，C' 与 Z(C')）是投射的（相应地，内射的）。证明从 H(C) 到 H(C') 的每个复形的态射都由从 C 到 C' 的某个态射诱导。
\item
  设 C 为 A-模复形；将 C 视为 A(ε)-模（X，第 27 页）。我们用
\end{enumerate}

\[
i: \mathbf {A} \rightarrow \mathbf {A} (\varepsilon)
\]

表示典范单射。

证明以下条件是等价的：

\begin{enumerate}
\def\labelenumi{(\roman{enumi})}
\item
  C 是 i-投射的（X，第 170 页，习题 19）。
\item
  它是 i-内射的（同上）。
\item
  复形 C 同伦于零。
\item
  存在一个分次 A-模 M 以及分次 A(ε)-模的同构 A(ε) ⊗A M → C。
\item
  存在一个分次 A-模 N 以及分次 A(ε)-模的同构 C → Hom\_A(A(ε), N)。
\end{enumerate}

\begin{enumerate}
\def\labelenumi{\arabic{enumi})}
\setcounter{enumi}{3}
\tightlist
\item
  \begin{enumerate}
  \def\labelenumii{\alph{enumii})}
  \tightlist
  \item
    设 C 为右有界 A-复形，且 C 与 H(C) 均为投射的。证明 C 是分裂的。
  \end{enumerate}
\end{enumerate}

\begin{enumerate}
\def\labelenumi{\alph{enumi})}
\setcounter{enumi}{1}
\tightlist
\item
  在环 B = A(ε)（X，第 27 页）上，我们考虑由 C\_n = B （对所有 n）且 d(b) = εb （对 b ∈ B）定义的复形 C。复形 C 是自由的且同调为零，但并非分裂的。
\end{enumerate}

\begin{enumerate}
\def\labelenumi{\arabic{enumi})}
\setcounter{enumi}{4}
\tightlist
\item
  设 C 为一个复形，并设 s 为 C 的一个分裂， \(\delta: \mathbf{C} \to \mathbf{C}\) 为一个次数为 -1 的分次同态，使得 \((d + \delta) \circ (d + \delta) = 0\) 。用 \(\mathbf{C}'\) 表示复形 \((\mathbf{C}, d + \delta)\) 。假设从 C 到 C 的 A-线性映射 \(\theta\) （定义为 \(\theta = 1_C + \delta s\) ）是双射。
\end{enumerate}

\begin{enumerate}
\def\labelenumi{\alph{enumi})}
\item
  为了使 \(s\theta^{-1}\) 成为 \(C'\) 的一个分裂，必要且充分条件是 \(\delta(Z(C)) \subset \theta(B(C))\) 。
\item
  假设 a) 的条件成立。证明存在直和分解： \(\mathbf{C} = \mathbf{B}(\mathbf{C}) \oplus \mathbf{Ker}(ds) = \mathbf{B}(\mathbf{C}') \oplus \mathbf{Ker}(ds)\) 。若 p（相应地， \(p'\) ）是像为 \(\mathbf{Ker}(ds)\) 且核为 \(\mathbf{B}(\mathbf{C})\) （相应地， \(\mathbf{B}(\mathbf{C}')\) ）的投影算子，则有 \(p' = p\theta^{-1}\) 。
\item
  我们进一步假设从 C 到 C 的 A-线性映射 \(\lambda\) （定义为 \(\lambda = 1_{C} + s\delta\) ）是双射。证明我们有 \(Z(C') = \lambda^{-1}(Z(C))\) 与 \(C = Z(C) \oplus \operatorname{Im}(sd) = Z(C') \oplus \operatorname{Im}(sd)\) 。若 \(q\) （相应地， \(q'\) ）是像为 \(Z(C)\) （相应地， \(Z(C')\) ）且核为 \(\operatorname{Im}(sd)\) 的投影算子，则有 \(q' = \lambda^{-1} q\) 。
\end{enumerate}

\begin{enumerate}
\def\labelenumi{\arabic{enumi})}
\setcounter{enumi}{5}
\tightlist
\item
  我们沿用习题 5 的记号；进一步假设 A 是交换的，复形 C 与 \(\mathbf{C}^{\prime}/\mathbf{B}(\mathbf{C}^{\prime})\) 是平坦的，并且存在 A 的一个理想 m，满足 \(\delta(\mathbf{C}) \subset \mathfrak{m}\mathbf{C}\) 与 \(\bigcap_{n} (\mathbf{B}(\mathbf{C}) + \mathfrak{m}^{n}\mathbf{C}) = \mathbf{B}(\mathbf{C})\) （若 A 是带极大理想 m 的局部诺特环且 \(C_{n}\) 对每个 n 都是有限生成的 A-模，则最后一个条件总满足（参见 AC，III，§ 3，n° 3，命题 6））。然后证明 \(s\theta^{-1}\) 是 \(C^{\prime}\) 的一个分裂。
\end{enumerate}

（平坦性假设蕴含 \(\operatorname{Im}(d + \delta) \cap \mathfrak{m}^n\mathbf{C} = (d + \delta)\, (\mathfrak{m}^n\,\mathbf{C})\) 对所有 \(n\) 成立；若 \(x \in \mathbf{Z}(\mathbf{C})\) ，通过归纳构造元素 \(y_n \in \mathbf{C}\) , \(z_n \in \mathbf{Z}(\mathbf{C}) \cap \mathfrak{m}^n\mathbf{C}\) ，使得 \(\delta(x) = \theta d(y_n) + \delta(z_n)\) 。利用习题 5 a) 得出结论。）

\begin{enumerate}
\def\labelenumi{\arabic{enumi})}
\setcounter{enumi}{6}
\tightlist
\item
  设 C 与 C' 是两个复形， \(s_{1}(\text{resp. } s')\) 为 C（相应地，C'）的一个分裂， \(u : C' \to C\) 为一个复形的态射。我们定义 Con(u) 的一个 A-自同态 \(\sigma\) ，其分次次数为 -1，通过设定
\end{enumerate}

\[
\sigma (y ^ {\prime}, x) = \left(- s ^ {\prime} (y ^ {\prime}), s (x) - s u s ^ {\prime} (y ^ {\prime})\right).
\]

证明为了使 \(\sigma\) 成为 Con \((u)\) 的一个分裂，必要且充分的条件是我们有 H(u) = 0。

\begin{enumerate}
\def\labelenumi{\arabic{enumi})}
\setcounter{enumi}{7}
\tightlist
\item
  设 \(0 \to M' \xrightarrow{\mu} M \xrightarrow{\nu} M'' \to 0\) 为一个 A-模的正合列，它可以被视为一个序列
\end{enumerate}

复形的正合列（在次数 ≠ 0 时为零）。证明复形的态射 φ : Con(u) → M''\,（X，第 39 页）是同伦等价，当且仅当正合列 0 → M' → M → M'' → 0 是分裂的。9) 设 u : C' → C 是 A-复形的一个态射。记 v : C → Coker(u) 为商映射。考虑分次 A-同态

\[
\alpha : (\operatorname{Ker} u) (- 1) \rightarrow \operatorname{Con} (u) \quad \text { and } \quad \varphi : \operatorname{Con} (u) \rightarrow \operatorname{Coker} (u)
\]

由 \(\alpha(x') = (x', 0)\) , \(\varphi(x', x) = v(x)\) （对 \(x' \in C'\) , \(x \bullet C\) ）定义。证明 \(\alpha, \varphi\) 是复形的态射，并且它们产生一个正合列

\[
\dots \rightarrow \mathrm{H} _ {n - 1} (\text { Ker } u) \xrightarrow {\mathrm{H} (\alpha)} \mathrm{H} _ {n} (\text { Con } (u)) \xrightarrow {\mathrm{H} (\varphi)} \mathrm{H} _ {n} (\text { Coker } u) \xrightarrow {\delta} \mathrm{H} _ {n - 2} (\text { Ker } u) \rightarrow \dots
\]

其中 δ 是相对于这两个正合列的连接同态的复合，

\[
0 \rightarrow \operatorname{Ker} u \rightarrow C ^ {\prime} \rightarrow \operatorname{Im} u \rightarrow 0 \quad \text { and } \quad 0 \rightarrow \operatorname{Im} u \rightarrow C \rightarrow \operatorname{Coker} u \rightarrow 0.
\]

\begin{enumerate}
\def\labelenumi{\arabic{enumi})}
\setcounter{enumi}{9}
\tightlist
\item
  设 \(u: \mathbf{C} \to \mathbf{C}'\) 为 A-复形的一个态射， \(\mathbf{P}\) 为一个投射 A-模， \(n\) 为一个整数 \(\geqslant 0\) ，
\end{enumerate}

\[
h: \mathrm{P} \rightarrow \mathrm{H} _ {n} (\text { Con } (u))
\]

一个 A-同态。证明存在一个复形 \(\tilde{C}\) 、复形的一个正合列

\[
0 \to \mathbf {C} \xrightarrow {i} \tilde {\mathbf {C}} \to \mathbf {P} (- n) \to 0
\]

以及一个态射 \(\tilde{u}:\tilde{\mathbf{C}}\to \mathbf{C}'\) ，使得 \(\tilde{u}\circ i = u\) ，并且态射 \(\mathbf{C}(i):\mathrm{Con}(u)\to \mathrm{Con}(\tilde{u})\) （由 \(i\) 诱导）满足以下性质：

α) 映射 \(\mathrm{H}_{p}(\mathrm{C}(i)) : \mathrm{H}_{p}(\mathrm{Con}(u)) \to \mathrm{H}_{p}(\mathrm{Con}(\tilde{u}))\) 对 \(p \neq n, n + 1\) 是双射的；

β) 存在一个正合列

\[
0 \rightarrow \mathrm{H} _ {n + 1} (\operatorname{Con} (u)) \xrightarrow {\mathrm{H} _ {n + 1} (\mathbf {C} (i))} \mathrm{H} _ {n + 1} (\operatorname{Con} (\tilde {u})) \rightarrow \mathrm{P} \xrightarrow {h} \mathrm{H} _ {n} (\operatorname{Con} (u)) \xrightarrow {\mathrm{H} _ {n} (\mathbf {C} (i))} \mathrm{H} _ {n} (\operatorname{Con} (\tilde {u})) \rightarrow 0.
\]

\begin{enumerate}
\def\labelenumi{\arabic{enumi})}
\setcounter{enumi}{10}
\tightlist
\item
  我们把 A-模的双复形称为一个三元组 (C, \(d'\) , \(d''\) )，其中 C 是一个类型为 \(\mathbf{Z} \times \mathbf{Z}\) 的 A-分次模， \(d': \mathbf{C} \to \mathbf{C}\) 是一个次数为 (-1, 0) 的分次 A-自同态，且 \(d'': \mathbf{C} \to \mathbf{C}\) 是一个次数为 (0, -1) 的分次 A-自同态，满足 \(d' \circ d' = d'' \circ d'' = d' \circ d'' + d'' \circ d' = 0\) 。我们还将使用由 \(\mathbf{C}^{p,q} = \mathbf{C}_{-p,-q}\) （对 \(p, q \in \mathbf{Z}\) ）定义的上升分次。
\end{enumerate}

我们用 \(Z'(C)\) , \(Z''(C)\) , \(B'(C)\) , \(B''(C)\) 表示子双复形 \(\text{Ker } (d')\) , \(\text{Ker } (d'')\) , \(\text{Im } (d')\) , \(\text{Im } (d'')\) ，并用 \(H'(C)\) （相应地， \(H''(C)\) ）表示双复形 \(Z'(C)/B'(C)\) （相应地， \(Z''(C)/B''(C)\) ）。

\begin{enumerate}
\def\labelenumi{\alph{enumi})}
\tightlist
\item
  我们设 \(d = d' + d''\) 。证明当 C 装备上与给定双分次相关联的全分次时（由 \(\mathbf{C}_n = \bigoplus_{p+q=n} \mathbf{C}_{p,q}\) 定义，参见 II，第 164 页），二元组 (C, d) 是一个复形。它称为
\end{enumerate}

与双复形 (C, \(d'\) , \(d''\) ) 相关联的复形；其同调记为 H(C)。

\begin{enumerate}
\def\labelenumi{\alph{enumi})}
\setcounter{enumi}{1}
\item
  设 (C, \(d'\) , \(d''\) ) 与 (K, \(\delta'\) , \(\delta''\) ) 是两个双复形。从 (C, \(d'\) , \(d''\) ) 到 (K, \(\delta'\) , \(\delta''\) ) 的一个态射是从 C 到 K 的一个次数为 (0, 0) 的分次 A-同态 \(u\) ，满足： \(\delta'\circ u = u \circ d'\) 与 \(\delta''\circ u = u \circ d''\) 。证明 \(u\) 诱导了相应复形之间的一个态射。
\item
  设 \(u, v\) 为从 C 到 K 的两个双复形的态射。称连接 \(u\) 与 \(v\) 的同伦为从 C 到 K 的一对 A-同态 \((s', s'')\) ，其次数分别为 (1, 0) 和 (0, 1)，使得
\end{enumerate}

\[
g - f = d ^ {\prime} \circ s ^ {\prime} + s ^ {\prime} \circ d ^ {\prime} + d ^ {\prime \prime} \circ s ^ {\prime \prime} + s ^ {\prime \prime} \circ d ^ {\prime \prime}; \quad s ^ {\prime} \circ d ^ {\prime \prime} + d ^ {\prime \prime} \circ s ^ {\prime} = 0; \quad s ^ {\prime \prime} \circ d ^ {\prime} + d ^ {\prime} \circ s ^ {\prime \prime} = 0.
\]

证明 \(s' + s''\) 是连接由 \(u\) 与 \(v\) 诱导的相应复形的态射的同伦。12) 设 \((\mathbf{C}, d', d'')\) 为一个双复形，使得 \(\mathrm{H}_{p,-p}''(\mathbf{C}) = 0\) 对所有 \(p \in \mathbf{Z}\) 成立， \(\mathrm{H}_{-q,q+1}'(\mathbf{C}) = 0\) 当 \(q \geqslant 0\) ， \(\mathrm{H}_{r,-r-1}'(\mathbf{C}) = 0\) 当 \(r \geqslant 0\) ，并且存在正整数 \(a\) 与 \(b\) 使得 \(\mathbf{C}_{a,-a} = \mathbf{C}_{-b,b} = 0\) 。证明 A-模 \(\mathrm{H}_{0,0}'(\mathbf{C})\) 为零。

\begin{enumerate}
\def\labelenumi{\arabic{enumi})}
\setcounter{enumi}{12}
\tightlist
\item
  设 I 为一个有限集；我们用 \((e_i)_{i\in I}\) 表示 Z-模 \(\mathbf{Z}^{(I)}\) 的典范基。我们称一个 A-模的 I-复形（相应地，I-预复形）为一个类型为 \(\mathbf{Z}^{I}\) 的分次 A-模 C，配备一族 A-自同态 \((d_i)_{i\in I}\) ，使得 \(d_i\) 的分次次数为 \((-e_i)\) ，满足：
\end{enumerate}

\[
d _ {i} \circ d _ {i} = 0 \quad \text { for every } \quad i \in I;
\]

\[
d _ {i} \circ d _ {j} + d _ {j} \circ d _ {i} = 0 (\text { resp. } d _ {i} \circ d _ {j} = d _ {j} \circ d _ {i}) \quad \text { for } \quad i, j \in I.
\]

若 Card (I) = 1，则一个 I-复形就是一个复形；若 Card (I) = 2，则一个 I-复形就是一个双复形（习题 11）。若 \((\mathbf{C}, (d_i))\) 与 \((\mathbf{C}', (d_i'))\) 是两个 I-复形（相应地，I-预复形），则从 \((\mathbf{C}, (d_i))\) 到 \((\mathbf{C}', (d_i'))\) 的一个态射是一个次数为零的分次同态 \(u: \mathbf{C} \to \mathbf{C}'\) ，使得 \(d_i' \circ u = u \circ d_i\) 对所有 \(i \in I\) 成立。

\begin{enumerate}
\def\labelenumi{\alph{enumi})}
\tightlist
\item
  设 (C, \((d_i)\) ) 为一个 I-复形。若赋予 C 由其分次导出的全分次（Z 型），证明 (C, \(\sum_{i\in I}d_i\) ) 是一个 A-模复形，称为与 I-复形 (C, \((d_i)\) ) 相关联的复形。
\end{enumerate}

证明 I-复形的态射诱导出相应复形的态射。

\begin{enumerate}
\def\labelenumi{\alph{enumi})}
\setcounter{enumi}{1}
\tightlist
\item
  设 C 是一个 I-预复形，并设 R 是 I 上的一个全序关系，记为 ≤。对于 \(k = (k_i)_{i \in I}, x \in C_k\) 与 \(i \in I\) ，我们设
\end{enumerate}

\(\delta_{i}(x) = (-1)^{\sum_{j < i} k_{j}} d_{i}(x)\) 。证明 \((C, (\delta_{i}))\) 是一个 I-复形，记为 \(C_{R}\) 。

\begin{enumerate}
\def\labelenumi{\alph{enumi})}
\setcounter{enumi}{2}
\item
  设 \(\mathbf{R}_{\alpha}, \mathbf{R}_{\beta}\) 为 I 上的两个全序关系，分别记为 \(\leqslant_{\alpha}\) 与 \(\leqslant_{\beta}\) 。设 \(u_{\beta \alpha}\) 为 C 的自同态，由 \(u(x) = (-1)^{\varepsilon(k)} x\) （对 \(x \in \mathbf{C}_k\) ，其中 \(\varepsilon(k) = \sum_{\substack{i < _{\alpha} j \\ j < _{\beta} i}} k_i k_j\) ）定义。证明 \(u_{\beta \alpha}\) 是从 \(\mathbf{C}_{\mathbf{R}_{\alpha}}\) 到 \(\mathbf{C}_{\mathbf{R}_{\beta}}\) 的 I-复形的态射。
\item
  设 S 为 I 上全体全序关系组成的集；我们在 S 上赋予平凡预序（即满足如下条件的预序： \(R_{\alpha} \leqslant R_{\beta}\) ，无论 \(R_{\alpha}, R_{\beta} \in S\) ）。证明系统 \((\mathbf{C}_{R\alpha}, u_{\alpha\beta})\) 是相对于 S 的 I-复形的投影系；其极限 C 是一个 I-复形，称为与 I-预复形 C 相关联的 I-复形。证明 I-预复形的每个态射都诱导出相关联 I-复形之间的一个态射。
\end{enumerate}

A-模的一个谱序列定义为 A-复形的一个序列 \((\mathrm{E}_r, d_r)_{r \geqslant 2}\) 以及一族映射 \(\alpha_r: \mathrm{H}(\mathrm{E}_r) \to \mathrm{E}_{r+1}\) ，使得：

\begin{enumerate}
\def\labelenumi{(\roman{enumi})}
\tightlist
\item
  复形 \(E_r\) 的分次是与一个双分次相关联的全分次
\end{enumerate}

\[
\mathrm{E} _ {r} = \bigoplus_ {p, q \in \mathbf {Z}} \mathrm{E} _ {r} ^ {p, q} (r \geqslant 2);
\]

我们有 \(d_r(E_r^{pq}) \subset E_r^{p + r,q - r + 1}\) 。

\begin{enumerate}
\def\labelenumi{(\roman{enumi})}
\setcounter{enumi}{1}
\tightlist
\item
  映射 \(\alpha_{r}\) 是双分次 A-模之间的同构。
\end{enumerate}

我们还将使用 \(E_{r}\) 的下降双分次：记 \(E_{pq}^{r} = E_{r}^{-p, -q}\) 。

\begin{enumerate}
\def\labelenumi{\alph{enumi})}
\tightlist
\item
  这里我们用 \(Z_{r+1}(E_r)\) （相应地， \(B_{r+1}(E_r)\) ）表示双分次 A-模 Ker \(d_r\) （相应地，Im \(d_r\) ）。设
\end{enumerate}

\[
p _ {r}: \mathrm{E} _ {r} \rightarrow \mathrm{E} _ {r} / \mathrm{B} _ {r + 1} (\mathrm{E} _ {r})
\]

商映射以及 \(i_{r}: \mathbf{E}_{r+1} \to \mathbf{E}_{r}/\mathbf{B}_{r+1}(\mathbf{E}_{r})\) 为 \(\alpha_{r}^{-1}\) 与自然单射的复合映射。对每个 \(r \geqslant 2\) ，我们通过对 \(k\) 的归纳，如下定义双分次子模 \(\mathbf{Z}_{r+k}(\mathbf{E}_{r})\) 与 \(\mathbf{B}_{r+k}(\mathbf{E}_{r})\) （它们是 \(\mathbf{E}_{r}\) 的子模）：

\[
\mathbf {Z} _ {r + k} (\mathrm{E} _ {r}) = p _ {r} ^ {- 1} \left(i _ {r} \left(\mathbf {Z} _ {r + k} (\mathrm{E} _ {r + 1})\right)\right); \quad \mathbf {B} _ {r + k} (\mathrm{E} _ {r}) = p _ {r} ^ {- 1} \left(i _ {r} \left(\mathbf {B} _ {r + k} (\mathrm{E} _ {r + 1})\right)\right).
\]

证明我们有以下包含关系：

\[
0 \subset \mathrm{B} _ {r + 1} (\mathrm{E} _ {r}) \subset \mathrm{B} _ {r + 2} (\mathrm{E} _ {r}) \subset \dots \subset \mathrm{Z} _ {r + 2} (\mathrm{E} _ {r}) \subset \mathrm{Z} _ {r + 1} (\mathrm{E} _ {r}) \subset \mathrm{E} _ {r}
\]

并且 \(\mathbf{E}_k\) 同构于 \(\mathbf{Z}_k(\mathbf{E}_r) / \mathbf{B}_k(\mathbf{E}_r)\) （对 \(k\geqslant r + 1\) ）。

我们设 \(Z_{\infty}(E_2) = \bigcap_{r \geqslant 3} Z_r(E_2), B_{\infty}(E_2) = \bigcup_{r \geqslant 3} B_r(E_2)\) 与 \(E_{\infty} = Z_{\infty}(E_2)/B_{\infty}(E_2)\) 。我们称谱序列是正则的，如果递减序列 \((Z_r(E_2^{pq}))_{r \geqslant 2}\) 对每一对 \((p, q)\) 都是平稳的。

\begin{enumerate}
\def\labelenumi{\alph{enumi})}
\setcounter{enumi}{1}
\item
  设 \(E = (E_r, d_r, \alpha_r)_{r \geqslant 2}\) 与 \(E' = (E_r', d_r', \alpha_r')_{r \geqslant 2}\) 为两个 A-模的谱序列；从 E 到 \(E'\) 的一个态射是一列 A-复形的态射 \(u_r : E_r \to E_r'\) ，它与双分次相容，使得 \(\alpha_r' \circ H(u_r) = u_{r+1} \circ \alpha_r\) 对所有 \(r \geqslant 2\) 成立。证明对于每一个 \(r \geqslant 2\) 以及每个 \(k \geqslant r + 1\) , \(u_r\) 诱导出从 \(Z_k(E_r)\) （相应地， \(B_k(E_r)\) ）到 \(Z_k(E_r')\) （相应地， \(B_k(E_r')\) ）的一个同态。特别地， \(u_2\) 诱导出同态 \(Z_\infty(u_2) : Z_\infty(E_2) \to Z_\infty(E_2')\) , \(B_\infty(u_2) : B_\infty(E_2) \to B_\infty(E_2')\) 与 \(u_\infty : E_\infty \to E_\infty'\) ；所有这些同态都与双分次相容。
\item
  我们假设 \(u_2\) 是一个同构。证明此时 \(u_r\) 对所有 \(r \geqslant 2\) 以及 \(\mathbf{B}_{\infty}(u_2)\) 都是同构。此外，若序列 E 和 E' 是正则的，则 \(\mathbf{Z}_{\infty}(u_2)\) 与 \(u_{\infty}\) 都是同构。
\item
  设 \(G = \bigoplus_{n \in Z} G^n\) 为一个分次 A-模，配备分次子 A-模的一个递减序列 \((F^p G)_{p \in \mathbb{Z}}\)
\end{enumerate}

使得 \(\bigcap_{p} F^{p} G = 0\) 且 \(\bigcup_{p} F^{p} G = G\) 。我们说谱序列 E 逼近 \((G, (F^{p} G)_{p \in \mathbb{Z}})\)

（或简称为 G），如果存在同构 \(\beta^{pq}: E_{\infty}^{pq} \to (F^p G)^{p+q}/(F^{p+1} G)^{p+q}\) 对所有

\[
(p, q) \in \mathbf {Z} \times \mathbf {Z}.
\]

称谱序列 E 收敛于 G，如果它逼近 G 并且此外：

\begin{enumerate}
\def\labelenumi{(\roman{enumi})}
\item
  谱序列 E 是正则的；
\item
  对于每个整数 n，存在一个整数 p，使得 \((\mathbf{F}^{p}\mathbf{G})^{n}=0\) 。
\end{enumerate}

假设 E（相应地 E′）收敛于 G（相应地 G′）。设 \(u: \mathrm{E} \to \mathrm{E}'\) 为谱序列的一个态射， \(v: \mathrm{G} \to \mathrm{G}'\) 为一个次数为 0 的分次同态，使得 \(v(\mathrm{F}^p \mathrm{G}) \subset \mathrm{F}^p \mathrm{G}'\) 对所有 \(p\) 成立，并且诱导的同态 \(\bar{v}^{p,n}: (\mathrm{F}^p \mathrm{G})^n / (\mathrm{F}^{p+1} \mathrm{G})^n \to (\mathrm{F}^p \mathrm{G}')^n / (\mathrm{F}^{p+1} \mathrm{G}')^n\) 满足 \(\bar{v}^{p,p+q} \circ \beta^{p,q} = \beta'^{p,q} \circ u_{\infty}^{p,q}\) 对每一对 \((p, q)\) 成立。证明如果 \(u_2^{pq}\) 对每一对 \((p, q)\) 都是同构的，则 \(v\) 是一个同构。

\begin{enumerate}
\def\labelenumi{\arabic{enumi})}
\setcounter{enumi}{14}
\tightlist
\item
  设 E 是一个 A-模的谱序列，它逼近一个分次 A-模 G（习题 14）。
\end{enumerate}

\begin{enumerate}
\def\labelenumi{\alph{enumi})}
\tightlist
\item
  假设 \(E_{2}^{p,q}=0\) 当 p\textless0。证明映射 \(\alpha_{r}\) 对每个 q 诱导出单射的 A-同态：
\end{enumerate}

\[
\mathrm{E} _ {\infty} ^ {0, q} \rightarrow \dots \rightarrow \mathrm{E} _ {3} ^ {0, q} \rightarrow \mathrm{E} _ {2} ^ {0, q}
\]

由此，通过与 \((\beta^{0q})^{-1}\) 复合以及与到商映射的复合，得到一个 A-同态 \(\gamma^q: G^q \to E_2^{0,q}\) 。

\begin{enumerate}
\def\labelenumi{\alph{enumi})}
\setcounter{enumi}{1}
\tightlist
\item
  假设 \(E_{2}^{pq} = 0\) 当 q \textless{} 0。证明映射 \(\alpha_{r}^{-1}\) 诱导出满射同态：
\end{enumerate}

\[
\mathrm{E} _ {2} ^ {p, 0} \rightarrow \mathrm{E} _ {3} ^ {p, 0} \rightarrow \dots \rightarrow \mathrm{E} _ {\infty} ^ {p, 0}
\]

由此，通过与 \(\beta^{p,0}\) 复合，得到一个 A-同态 \(\delta^p:E_2^{p,0}\to G^p\) 。c) 我们假设 a) 和 b) 的假设均成立，即 \(E_{2}^{pq} = 0\) （若 p \textless{} 0 或 q \textless{} 0）。证明序列

\[
0 \rightarrow \mathrm{E} _ {2} ^ {1, 0} \xrightarrow {\delta^ {1}} \mathrm{G} ^ {1} \xrightarrow {\gamma^ {1}} \mathrm{E} _ {2} ^ {0, 1} \xrightarrow {d _ {2}} \mathrm{E} _ {2} ^ {2, 0} \xrightarrow {\delta^ {2}} \mathrm{G} ^ {2}
\]

是正合的。

\begin{enumerate}
\def\labelenumi{\alph{enumi})}
\setcounter{enumi}{3}
\item
  假设 \(\mathbf{E}_2^{pq} = 0\) （若 \(p > 0\) 或 \(q > 0\) ）；定义一个与 \(c\) ) 中类似的正合列。
\item
  假设存在整数 \(r, d\) ，满足 \(d \geqslant r \geqslant 2\) ，使得 \(\mathrm{E}_{r}^{pq} = 0\) （若 \(p \neq 0\) , \(d\) ）。定义一个正合列：
\end{enumerate}

\[
\dots \rightarrow \mathrm{E} _ {r} ^ {d, n - d} \rightarrow \mathrm{G} ^ {n} \rightarrow \mathrm{E} _ {r} ^ {0, n} \rightarrow \mathrm{E} _ {r} ^ {d, n + 1 - d} \rightarrow \mathrm{G} ^ {n + 1} \rightarrow \dots .
\]

若同样 \(E_{r}^{pq} = 0\) （当 \(q \neq 0\) , d），证明存在正合列：

\[
\dots \rightarrow \mathrm{E} _ {r} ^ {n, 0} \rightarrow \mathrm{G} ^ {n} \rightarrow \mathrm{E} _ {r} ^ {n - d, d} \rightarrow \mathrm{E} _ {r} ^ {n + 1, 0} \rightarrow \dots .
\]

\begin{enumerate}
\def\labelenumi{\alph{enumi})}
\setcounter{enumi}{5}
\tightlist
\item
  假设 \(\mathbf{E}_2^{pq} = 0\) （若 \(q \neq 0\) ）（相应地， \(p \neq 0\) ）。证明同态
\end{enumerate}

\[
\delta^ {n}: \mathrm{E} _ {2} ^ {n, 0} \rightarrow \mathrm{G} ^ {n} (\text { resp. } \gamma^ {n}: \mathrm{G} ^ {n} \rightarrow \mathrm{E} _ {2} ^ {0, n})
\]

对每个 n 都是双射。

¶ 16) 设 C 是一个 A-复形， \((F^{p}C)_{p\in\mathbf{Z}}\) 是 C 的子复形的一个递减序列。

\begin{enumerate}
\def\labelenumi{\alph{enumi})}
\tightlist
\item
  对于 \(p,q,r\in\mathbf{Z}\) ，设 \(M_{r}^{pq}\) 为 \((F^{p}C)^{p+q}\) 中满足 \(dx\in F^{p+r}C\) 的元素 x 组成的集合；对 \(r\geqslant0\) ，我们设 \(E_{r}^{pq}=M_{r}^{pq}/(d(M_{r-1}^{p-r+1,q+r-2})+M_{r}^{p+1,q-1})\) ，并用 \(d_{r}:E_{r}^{pq}\to E_{r}^{p+r,q-r+1}\) 表示由 \(d\) 通过取商得到的同态。证明对 \(r\geqslant0\) 存在同构 \(\alpha_{r}:H(E_{r})\to E_{r+1}\) ，使得序列 \((E_{r},d_{r})_{r\geqslant2}\) （配备映射 \(\alpha_{r}\) ）是一个谱序列。我们有
\end{enumerate}

\[
\mathrm{E} _ {0} ^ {p q} = (\mathrm{F} ^ {p} \mathrm{C}) ^ {p + q} / (\mathrm{F} ^ {p + 1} \mathrm{C}) ^ {p + q}.
\]

\begin{enumerate}
\def\labelenumi{\alph{enumi})}
\setcounter{enumi}{1}
\item
  我们假设 \(\bigcup_{p} \mathbf{F}^{p} \mathbf{C} = \mathbf{C}\) 并且对于每个 \(n\) ，存在一个整数 \(p\) 使得 \((\mathbf{F}^{p} \mathbf{C})^{n} = 0\) 。证明在 \(a)\) 中定义的谱序列收敛（习题 14）于分次 A-模 \(\mathrm{H}(\mathbf{C})\) ，它配备分次子模序列 \(\mathbf{F}^{p} \mathbf{H}(\mathbf{C}) = \operatorname{Im} \mathbf{H}(i_{p})\) ，其中 \(i_{p}: \mathbf{F}^{p} \mathbf{C} \to \mathbf{C}\) 是典范单射。
\item
  若 \(\mathbf{F}^0\bar{\mathbf{C}} = \mathbf{C}\) ，则 A-模 \(\mathbf{E}_r^{pq}\) 对 \(p < 0\) 为零；若 \((\mathbf{F}^p\mathbf{C})^n = 0\) （当 \(p > n\) ），则有 \(\mathbf{E}_r^{pq} = 0\) （对 \(q < 0\) ）。d) 设 \(\mathbf{C}'\) 为第二个复形， \((\mathbf{F}^p\mathbf{C}')_{p\in \mathbb{Z}}\) 为 \(\mathbf{C}'\) 的子复形的一个递减序列， \(\mathbf{E}'\) 为相应的谱序列。设 \(u:\mathbf{C}\to \mathbf{C}'\) 是复形的态射，使得 \(u(\mathbf{F}^p\mathbf{C})\subset \mathbf{F}^p\mathbf{C}'\) 对所有 \(p\) 成立。证明 \(u\) 诱导谱序列的一个态射 \(\tilde{u}:\mathbf{E}\to \mathbf{E}'\) 。
\item
  沿用 \(d\) ) 的记号，设 \(v: \mathbf{C} \to \mathbf{C}'\) 是第二个态射，使得 \(v(\mathbf{F}^p \mathbf{C}) \subset \mathbf{F}^p \mathbf{C}'\) 对所有 \(p\) 成立，并设 \(h\) 为联系 \(u\) 与 \(v\) 的一个同伦。设 \(k\) 为一个整数 \(\geqslant 0\) ，使得 \(h(\mathbf{F}^p \mathbf{C}) \subset \mathbf{F}^{p-k} \mathbf{C}'\) 对所有 \(p\) 成立。证明 \(\tilde{u}_r = \tilde{v}_r\) 当 \(r > k\) 。
\end{enumerate}

\begin{enumerate}
\def\labelenumi{\arabic{enumi})}
\setcounter{enumi}{16}
\tightlist
\item
  设 C 是一个双复形（X，第 174 页，习题 11）。
\end{enumerate}

\begin{enumerate}
\def\labelenumi{\alph{enumi})}
\tightlist
\item
  我们定义与 C 相关联的复形的子复形的两个递减序列 \(({}^{\prime}\mathbf{F}^p\mathbf{C})_{p\in \mathbf{Z}}\) 与 \(({}^{\prime\prime}\mathbf{F}^q\mathbf{C})_{q\in \mathbf{Z}}\) ，通过设定
\end{enumerate}

\[
\left(^ {\prime} \mathrm{F} ^ {p} \mathrm{C}\right) ^ {n} = \bigoplus_ {i \geqslant p} \mathrm{C} ^ {i, n - i} \quad \left(^ {\prime \prime} \mathrm{F} ^ {q} \mathrm{C}\right) ^ {n} = \bigoplus_ {j \geqslant q} \mathrm{C} ^ {n - j, j}.
\]

设′E 与″E 为相应的谱序列（习题 16，a））。证明对每一对 \((p, q)\) ，A-模 \({}^{\prime}E_2^{pq}\) （相应地， \({}^{\prime\prime}E_2^{pq}\) ）同构于 \(H^{\prime\prime p}(H^{\prime q}(C))\) （相应地， \(H^{\prime p}(H^{\prime\prime q}(C))\) ）。

\begin{enumerate}
\def\labelenumi{\alph{enumi})}
\setcounter{enumi}{1}
\item
  假设存在一个整数 \(n\) 使得 \(C^{pq} = 0\) （当 \(p \geqslant n\) 或 \(q \leqslant n\) ）。证明谱序列 ′E 收敛于 H(C)。如果同样存在一个整数 \(m\) 使得 \(C^{pq} = 0\) （当 \(p \leqslant m\) 或 \(q \geqslant m\) ），则谱序列 ″E 收敛于 H(C)。
\item
  我们假设 \(C^{pq} = 0\) （当 p \textless{} 0 或 q \textless{} 0）。证明同态 \(\delta^{p}: E_{2}^{p0} \to H^{p}(C)\) （习题 15）在取同调后，由复形的典范单射 ``Z\^{}\{0\}(C) → C'' 诱导而来。同样地，同态 ``E\_\{2\}\^{}\{p0\} → H\^{}\{p\}(C)'' 来自单射 ``Z\^{}\{0\}(C) → C''。
\item
  设 \(\tilde{C}\) 为第二个双复形，′E 与 ″E 为两个相应的谱序列。证明双复形的每个态射 \(u: C \to \tilde{C}\) 诱导谱序列的态射 ′u : ′E → ′E 与 ″u : ″E → ″E。证明如果两个态射 \(u\) , \(v\) （从 \(C\) 到 \(\tilde{C}\) ）是同伦的，则相应的态射 ′u 与 ′v（相应地，″u 与 ″v）相等（利用习题 16，e））。
\end{enumerate}

\begin{enumerate}
\def\labelenumi{\arabic{enumi})}
\setcounter{enumi}{17}
\tightlist
\item
  设 E 是一个谱序列（习题 14），使得 A-模 \(\mathbf{E}_2\) 具有有限长度。设 \(\mathcal{C}\) 是有限长度 A-模的同构类组成的集合；证明有 \(\chi_{\mathcal{G}}(\mathbf{E}_2) = \chi_{\mathcal{G}}(\mathbf{E}_r) = \chi_{\mathcal{G}}(\mathbf{E}_{\infty})\) 对所有 \(r \geqslant 2\) 成立（X，第 41 页）。如果此外 E 收敛于一个分次 A-模 G，则后者具有有限长度，并且有 \(\chi_{\mathcal{G}}(\mathbf{G}) = \chi_{\mathcal{G}}(\mathbf{E}_2)\) 。
\end{enumerate}

¶ 19) 我们称一个单纯复形为一个对 (K, F)，其中 K 是一个集合，F 是 K 的有限子集组成的一个集合，称为 K 的面，使得一个面的每个子集都是一个面，且 K 的每个点至少属于一个面。从一个单纯复形 (K, F) 到一个单纯复形 (L, G) 的单纯映射是从 K 到 L 的一个映射，它把 K 的每个面映到 L 的一个面。a) 设 n 是一个整数 ≥ 0。我们称单纯复形 K 的一个 n-单形为 K 的元素的序列 (α₀, \ldots, αₙ)，使得诸 αᵢ 属于 K 的同一个面。我们用 Sₙ(K, A) 表示由 n-单形（n ≥ 0）生成的自由 A-模，并对 n \textless{} 0 设 Sₙ(K, A) = 0。我们如下定义一个 A-同态 \(d_{n}:\mathbf{S}_{n}(\mathbf{K},\mathbf{A})\to \mathbf{S}_{n - 1}(\mathbf{K},\mathbf{A})\) ：设定 \(d_{n} = 0\) 当 \(n\leqslant 0\) ，且 \(d_{n}(\alpha_{0},\dots,\alpha_{n}) = \sum_{i = 0}^{n}(-1)^{i}(\alpha_{0},\dots,\theta_{i},\dots,\alpha_{n})\)

对所有 \(n\) -单形 \((\alpha_0, \ldots, \alpha_n)\) （字母上方的符号表示该字母应被省略）。证明 \(d_{n-1} \circ d_n = 0\) 对所有 \(n\) 成立，使得 \((S(K, A), d)\) 是 A-模的一个复形。以同样的方式定义 A-复形 \(C(K, A)\) ，通过设定 \(C^n(K, A) = \text{Hom}_Z(S_n(K, Z), A)\) , \(\delta^n = \text{Hom}(d_n, 1_A)\) （对所有 \(n\) ）。对 \(n \geqslant 0\) ，有时用 \(H_n(K, A)\) （相应地， \(H^n(K, A)\) ）表示 A-模 \(H_n(S(K, A))\) （相应地， \(H^n(C(K, A)))\) 。 b) 设 K, L 为两个单纯复形， \(f: K \to L\) 为一个单纯映射。证明 \(f\) 诱导复形的态射 \(S(f): S(K, A) \to S(L, A)\) 与 \(C(f): C(L, A) \to C(K, A)\) ，从而诱导 A-同态 \(f_*: H(K, A) \to H(L, A)\) 与 \(f^*: H'(L, A) \to H'(K, A)\) 。

\begin{enumerate}
\def\labelenumi{\alph{enumi})}
\setcounter{enumi}{2}
\tightlist
\item
  设 \(g: \mathbf{K} \to \mathbf{L}\) 为第二个单纯映射。我们说 \(f\) 与 \(g\) 是单纯同伦的，如果对每个面 \(\sigma\) （属于 \(\mathbf{K}\) ），集合 \(f(\sigma) \cup g(\sigma)\) 是 \(\mathbf{L}\) 的一个面。证明态射 \(\mathbf{S}(f)\) 与 \(\mathbf{S}(g)\) （相应地， \(\mathbf{C}(f)\) 与 \(\mathbf{C}(g)\) ）此时是同伦的（通过设定 \(h\) 来定义一个同伦
\end{enumerate}

\[
h (\alpha_ {0}, \dots , \alpha_ {n}) = \sum_ {i = 0} ^ {n} (- 1) ^ {i} (f (\alpha_ {0}), \dots , f (\alpha_ {i}), g (\alpha_ {i}), \dots , g (\alpha_ {n}))
\]

对每个 n-单形 \((\alpha_{0}, \ldots, \alpha_{n})\) 。

\begin{enumerate}
\def\labelenumi{\alph{enumi})}
\setcounter{enumi}{3}
\tightlist
\item
我们称单纯复形 K 为锥形的，如果存在 K 中的一点 s，使得对 K 的每个面 σ，σ ∪ \{s\} 是一个面。证明此时存在从 S(K, A) 和从 C(K, A) 到 A 上的同伦等价。e) 设 D(K, A) 表示 S(K, A) 的子复形，它由使得两个 αᵢ 相等的单形 (α₀, \ldots, αₙ) 以及形如 (α₀, \ldots, αₙ) - ε(σ). (ασ(0), \ldots, ασ(n)) 的元素（对每个 n-单形 (α₀, \ldots, αₙ) 以及 \{0, \ldots, n\} 的每个置换 σ）生成。证明 D(K, A) 是 S(K, A) 的一个子复形，并且商映射 p : S(K, A) → S(K, A)/D(K, A) 是同伦等价（通过在 K 上选取一个全序，并把单形 (α₀, \ldots..., αₙ) 模 D(K, A) 的类对应到单形 ε(σ). (ασ(0), \ldots, ασ(n))—其中 σ 是满足 ασ(0) \textless{} \ldots{} \textless{} ασ(n) 的置换，若诸 αᵢ 互不相同；否则对应到 0—即可得到一个同伦逆）。
\end{enumerate}

*20) 设 R 为一个完备离散赋值环（AC，VI，§ 3，第 6 节），其剩余域 k 的特征为 p \textgreater{} 0。假设 R 的分式域的特征为 0。设 v 表示 R 的规范化赋值，m 为 R 的极大理想，R* 为 R 的单位乘法群。对每个整数 \(n \geqslant 0\) ，设 R*(n) 为 R* 中由满足 \(v(x - 1) \geqslant n\) 的元素 x 构成的子群。 a) 设 r 为整数， \(r > v(p)/(p - 1)\) 。证明级数

\[
\exp (x) = \sum_ {i = 0} ^ {+ \infty} x ^ {i} / i!
\]

对 \(x \in \mathfrak{m}^r\) 收敛，并定义了从加法群 \(\mathfrak{m}^r\) 到群 \(\mathbf{R}^*(r)\) 上的一个同构。 b) 假设 \(k\) 是一个具有 \(q\) 个元素的有限域。设 \(n\) 为一个整数 \(\geqslant 1\) , \(\mu_n\) 为元素 \(x\) （ \(\mathbf{R}^*\) 中）满足 \(x^n = 1\) 所构成的子群。证明群 \(\mu_n\) 与 \(\mathbf{R}^*/(\mathbf{R}^*)^n\) 都是有限的，并且有

\[
\operatorname{Card} \left(\mathrm{R} ^ {*} / \mathrm{R} ^ {* n}\right) / \operatorname{Card} \left(\mu_ {n}\right) = q v (n).
\]

（设 \(\mathcal{C}\) 为有限 \(\mathbf{Z}\) -模的类组成的集合， \(\varphi : \mathrm{K}(\mathcal{C}) \to \mathbf{Q}^*\) 为如下定义的同态

\[
\varphi ([ \mathrm{G} ]) = \text { Card } (\mathrm{G}) \quad \text { for } \quad \mathrm{G} \in \mathscr {C},
\]

\(C^{(r)}\) 为满足如下条件的 Z-模复形： \(C_{p}^{(r)} = 0\) （当 \(p \neq 0, 1\) ）, \(C_{0}^{(r)} = C_{1}^{(r)} = R^{*}(r)\) , \(d_{1}(x) = x^{n}\) （当 \(x \in R^{*}(r)\) ）。证明 \(\varphi(H(C^{(r)}))\) 与 r 无关，然后使用 a)。

\begin{enumerate}
\def\labelenumi{\arabic{enumi})}
\setcounter{enumi}{20}
\tightlist
\item
  假设环 A 是交换环 k 上的一个交换代数；A 的 k-导子组成的 A-模 \(D_{k}(A)\) 于是被等同于 A-模 \(\Omega_{A/k}^{1}\) 的对偶，这在 \(\Omega_{A/k}\) 上定义了一个 \(\Lambda_{\mathbf{A}}(\mathbf{D}_{k}(\mathbf{A}))\) 左模结构（参见 III，第 165 页）。证明公式
\end{enumerate}

\[
\begin{array}{l} \left(\mathbf {X} _ {0} \wedge \dots \wedge \mathbf {X} _ {p}\right) \rfloor d \omega = \sum_ {i = 0} ^ {p} (- 1) ^ {i} \mathbf {X} _ {i} \left(\left(\mathbf {X} _ {0} \wedge \dots \wedge \hat {\mathbf {X}} _ {i} \wedge \dots \wedge \mathbf {X} _ {p}\right) \rfloor \omega\right) \\ \quad + \sum_ {0 \leqslant i <   j \leqslant p} (- 1) ^ {i + j} \left([ \mathbf {X} _ {i}, \mathbf {X} _ {j} ] \wedge \mathbf {X} _ {0} \wedge \dots \wedge \hat {\mathbf {X}} _ {i} \wedge \dots \wedge \hat {\mathbf {X}} _ {r} \wedge \dots \wedge \mathbf {X} _ {p}\right) \rfloor d \omega . \end{array}
\]

其中 \(\omega\in\Omega_{\mathbf{A}/k}^{p},\mathbf{X}_{0},\ldots,\mathbf{X}_{p}\in\mathbf{D}_{k}(\mathbf{A})\) ，而字母上的符号 \(^{\wedge}\) 表示该字母应被省略。（化归到 \(p=1\) 的情形。）

\begin{enumerate}
\def\labelenumi{\arabic{enumi})}
\setcounter{enumi}{21}
\tightlist
\item
  我们沿用前一习题的约定；若 M 是一个 A-模， \(\nabla\) 为 M 上的一个联络且若 \(X \in D_k(A)\) ，我们用 \(\nabla_X\) 表示 M 的 \(k\) -自同态 \((1_M \otimes X) \circ \nabla\) 。
\end{enumerate}

\begin{enumerate}
\def\labelenumi{\alph{enumi})}
\item
  设 \(\mathbf{M}_1, \mathbf{M}_2\) 为两个 A-模，分别配备联络 \(^1\nabla\) , \(^2\nabla\) ；证明可以如下定义 \(\nabla\) （ \(\mathbf{M}_1 \oplus \mathbf{M}_2\) （相应地， \(\mathbf{M}_1 \otimes_{\mathbf{A}} \mathbf{M}_2\) ，相应地， \(\operatorname{Hom}_{\mathbf{A}}(\mathbf{M}_1, \mathbf{M}_2)\) ）上的一个联络），通过设定： \(\nabla(m_1 + m_2) = ^1\nabla(m_1) + ^2\nabla(m_2)\) （当 \(m_1 \in \mathbf{M}_1\) , \(m_2 \in \mathbf{M}_2\) ）（相应地， \(\nabla_{\mathbf{X}}(m_1 \otimes m_2)) = ^1\nabla_{\mathbf{X}}(m_1) \otimes m_2 + m_1 \otimes ^2\nabla_{\mathbf{X}}(m_2)\) 当 \(\mathbf{X} \in \mathbf{D}_k(\mathbf{A})\) , \(m_1 \in \mathbf{M}_1\) 与 \(m_2 \in \mathbf{M}_2\) ，相应地， \(\nabla_{\mathbf{X}}(u) = ^2\nabla_{\mathbf{X}} \circ u - u \circ ^1\nabla_{\mathbf{X}}\) 当 \(\mathbf{X} \in \mathbf{D}_k(\mathbf{A})\) , \(u \in \operatorname{Hom}_{\mathbf{A}}(\mathbf{M}_1, \mathbf{M}_2)\) 。
\item
  设 M 为一个 A-模， \(\nabla: \mathbf{M} \to \mathbf{M} \otimes_{\mathbf{A}} \Omega_{\mathbf{A}/k}^{1}\) 为 M 上的一个联络。若 X, Y ∈ D \(_{k}\) (A)，我们用 R(X, Y) 表示 M 的如下 A-自同态：R(X, Y)(m) = \textless R(m), X ∧ Y\textgreater{} 对所有 m ∈ M 成立。证明公式：R(X, Y) = {[}∇X, ∇Y{]} - ∇{[}X,Y{]} （利用习题 21）。
\item
  假设 A-模 M 是投射的，从而曲率同态 R 被等同于 \(\operatorname{End}_{\mathbf{A}}(\mathbf{M}) \otimes_{\mathbf{A}} \Omega_{\mathbf{A}/k}^{2}\) 的一个元素。若给 \(\operatorname{End}_{\mathbf{A}}(\mathbf{M})\) 配备 \(\nabla\) （在 \(a\) ) 中定义的联络），证明有 \(\nabla^{2}\mathbf{R} = 0\) 。
\end{enumerate}

